% ---- pg-549 （书页 537） ----

\begin{mvtable}
\mvtabcap{续表}
\footnotesize
\begin{tabular}{lcccccccc}
\toprule
& \multicolumn{8}{c}{变\quad 量}\\
\cmidrule(lr){2-9}
公\quad 司 & $X_1$ & $X_2$ & $X_3$ & $X_4$ & $X_5$ & $X_6$ & $X_7$ & $X_8$\\
\midrule
17. 圣地亚哥天然气与电力公司 & 0.76 & 6.4 & 136.0 & 61.9 & 9.0 & $5\,714.0$ & 8.3 & 1.920\\
18. 南方公司 & 1.05 & 12.6 & 150. & 56.7 & 2.7 & $10\,140.0$ & 0 & 1.108\\
19. 得克萨斯公用事业公司 & 1.16 & 11.7 & 104.0 & 54.0 & $-2.1$ & $13\,507.0$ & 0 & 0.636\\
20. 威斯康星电力公司 & 1.20 & 11.8 & 148.0 & 59.9 & 3.5 & $7\,287.0$ & 41.1 & 0.702\\
21. 联合照明公司 & 1.04 & 8.6 & 204. & 61.0 & 3.5 & $6\,650.0$ & 0 & 2.116\\
22. 弗吉尼亚电力与电力公司 & 1.07 & 9.3 & 174.0 & 54.3 & 5.9 & $10\,093.0$ & 26.6 & 1.306\\
\bottomrule
\end{tabular}

\vspace{0.4em}
\begin{flushleft}
\noindent 变量定义:\quad $X_1$: 固定费用周转比(收入/债务)\\
\hspace*{5.2em}$X_2$: 资本回报率\\
\hspace*{5.2em}$X_3$: 每千瓦容量成本\\
\hspace*{5.2em}$X_4$: 年载荷因子\\
\hspace*{5.2em}$X_5$: 自 1974 年至 1975 年高峰期千瓦时需求增长\\
\hspace*{5.2em}$X_6$: 销售量(年千瓦时用量)\\
\hspace*{5.2em}$X_7$: 核能所占百分比\\
\hspace*{5.2em}$X_8$: 总燃料成本(美分/千瓦时)
\end{flushleft}

\vspace{0.3em}
\noindent 资料来源: 数据经 H. E. Thomson 许可使用.
\end{mvtable}

\noindent 我们用积矩相关系数来量度各对变量之间的相似性. 相关系数矩阵在\hyperref[tab:12.5]{表 12.5} 中给出.

\begin{mvtable}
\phantomsection\label{tab:12.5}\mvtabcap{表 12.5\quad 各对变量之间的相关系数(公用事业公司数据)}

\noindent\rule{\linewidth}{0.4pt}

\vspace{0.4em}
\[
\begin{array}{cccccccc}
X_1 & X_2 & X_3 & X_4 & X_5 & X_6 & X_7 & X_8
\end{array}
\]

\vspace{-1.0em}
\[
\begin{bmatrix}
1.000 & & & & & & & \\
0.643 & 1.000 & & & & & & \\
-0.103 & -0.348 & 1.000 & & & & &\\
-0.082 & -0.086 & 0.100 & 1.000 & & & &\\
-0.259 & -0.260 & 0.435 & 0.034 & 1.000 & & &\\
-0.152 & -0.010 & 0.028 & -0.288 & 0.176 & 1.000 & &\\
0.045 & 0.211 & 0.115 & -0.164 & -0.019 & -0.374 & 1.000 &\\
-0.013 & -0.328 & 0.005 & 0.486 & -0.007 & -0.561 & -0.185 & 1.000
\end{bmatrix}
\]
\end{mvtable}

\noindent 当样本相关系数被用作相似性量度时, 具有很大负相关系数的变量对被认为极不相似.
具有很大正相关系数的变量对则被认为非常相似. 在这种情形下, 两个聚类之间的“距离”用这两个聚类的对象之间的最小相似性来量度.
将完全连接算法用于上述相似性矩阵, 便产生\hyperref[fig:12.8]{图 12.8} 中的树图.

我们看到变量 1 和 2(固定费用周转比和资本回报率), 变量 4 和 8(年载荷因子和总燃料成本), 以及变量 3 和 5(每千瓦容量成本和高峰期千瓦时需求增长)在中等相似性水平上形成聚类.
变量 7(核能所占百分比)及 6(销售量)在最后一步之前始终自成一聚类. 最后的合并将聚类 $(12478)$ 和聚类 $(356)$ 带到一起. \hfill$\blacksquare$

\noindent 同单连接法的情形一样, 对于完全连接聚类, 距离(相似性)的保序变换不影响聚类的结果.

\mvsub{平均连接}

平均连接法将两个聚类之间的距离作为两个聚类间所有各对项目之间的平均距离来处
